\input{../tex-inputs/latex-leseansicht-vorspann}

\section[Theodor von Sosnosky an Arthur Schnitzler, 27. 12. 1899]{L04225 Theodor von Sosnosky an Arthur Schnitzler, 27. 12. 1899}
\nopagebreak\mylabel{L04225v}
\rehead{ }\normalsize\beginnumbering\briefempfaengerindex{Schnitzler, Arthur@\textsc{Schnitzler, Arthur}!zzzSosnosky, Theodor von@\emph{von Theodor von Sosnosky}!1899-12-271@{27. 12. 1899}|(be}
\toendnotes[C]{\smallbreak\pagebreak[2]}
\correspDesc{Versand  durch Theodor von Sosnosky am 27. 12. 1899 in Kremsmünster
\newline{}Erhalt  durch Arthur Schnitzler im Zeitraum [28. 12. 1899 – 1. 1. 1900?] in Wien}\toendnotes[C]{\smallbreak}
\Standort{DLA, A:Schnitzler, HS.1985.1.4640.}
\physDesc{Brief, 1 Blatt, 3 Seiten, 1946 Zeichen
\newline{}Handschrift: schwarze Tinte, deutsche Kurrent
\newline{}Schnitzler: mit Bleistift beschriftet: »Sosnosky« und datiert: »27/12 99« }\toendnotes[C]{\smallbreak}
\pstart{}{\pb}Sehr geehrter Herr Doktor!\pend\vspace{0.5em}
\pstart
           Im Begriffe eine Anthologie »\textsc{Die deutsche Lyrik des 19. Jahrhunderts\pwindex{deutsche Lyrik des 19. Jahrhunderts. Eine poetische Revue@\emph{Die deutsche Lyrik des 19. Jahrhunderts. Eine poetische Revue}|pw}}«
      zusa{\geminationm}enzuſtellen, worin ich ein vollſtändiges Spiegelbild der deutschen Lyrik
      von \textsc{Körner\pwindex{Körner, Theodor 23.\,9.\,1791 Dresden – 26.\,8.\,1813@\textsc{Körner, Theodor} (23.\,9.\,1791 Dresden – 26.\,8.\,1813), \emph{Schriftsteller, Dramatiker, Lyriker}|pw}} bis zu \textsc{Hoffmannsthal\pwindex{Hofmannsthal, Hugo von 1.\,2.\,1874 Wien – 15.\,7.\,1929 Rodaun@\textsc{Hofmannsthal, Hugo von} (1.\,2.\,1874 Wien – 15.\,7.\,1929 Rodaun), \emph{Schriftsteller}|pw}} und
      \textsc{Stefan George\pwindex{George, Stefan 17.\,7.\,1868 Büdesheim – 4.\,12.\,1933 Minusio@\textsc{George, Stefan} (17.\,7.\,1868 Büdesheim – 4.\,12.\,1933 Minusio), \emph{Schriftsteller, Übersetzer}|pw}} geben will, alſo keine konventionelle Blütenleſe für das deutſche
      Haus, geſtatte ich mir, bei Ihnen
      anzufragen, ob Sie es gerne ſähen,
      wenn auch lyrische Gedichte von Ihnen
      in die Sa{\geminationm}lung\pwindex{deutsche Lyrik des 19. Jahrhunderts. Eine poetische Revue@\emph{Die deutsche Lyrik des 19. Jahrhunderts. Eine poetische Revue}|pwv} aufgeno{\geminationm}en würden?
      Mir wäre es ein Vergnügen, Ihren
      Namen darin vertreten zu ſehen und
      Sie dem Publikum, das Sie bisher nur
      als Dramatiker und Novelliſt kennt
      und ſchätzt, nun auch als Lyriker
                {\pb}vorzuführen. Obwohl ich noch nie ein
      lyrisches Gedicht aus Ihrer Feder geleſen
      habe, so habe ich in Ihren Arbeiten einen
      ausgeſprochen lyrischen Zug zu erkennen
      geglaubt, der mich annehmen ließ, daß
      Sie auch lyriſch thätig ſind, eine Annahme, die Herr G. Schwarzkopf\pwindex{Schwarzkopf, Gustav 7.\,11.\,1853 Wien – 13.\,11.\,1939 ebd.@\textsc{Schwarzkopf, Gustav} (7.\,11.\,1853 Wien – 13.\,11.\,1939 ebd.), \emph{Schriftsteller}|pw} auf meine
      Frage beſtätigte. Dies der Grund, warum
      ich mich an Sie wende, Herr Doktor!\pend
           
\pstart
           Sollten Sie auf meinen Vorschlag
      eingehen, ſo würde ich bitten, mir eine
      Auswahl solcher Gedichte zu ſenden, die
      Sie für Ihre gelungenſten und für
      Ihre Eigenart meiſt charakteriſtischen
      halten; aus dieſen würde ich dann
      die mir am meiſten zuſagenden 2–5
      auswählen {\kaufmannsund} kopiren, worauf ich Ihnen
         Alles wieder zurückſenden möchte.\pend
           
\pstart
           Die Anthologie\pwindex{deutsche Lyrik des 19. Jahrhunderts. Eine poetische Revue@\emph{Die deutsche Lyrik des 19. Jahrhunderts. Eine poetische Revue}|pwv} ſoll bei Cotta\orgindex{J.G. Cotta’sche Buchhandlung Nachfolger@J.G. Cotta’sche Buchhandlung Nachfolger|pw} erſcheinen, doch ist es noch nicht sicher, denn
      obgleich er sich für sie intereſſirt
      und dies durch Zuſendung von Material
                {\pb}auch bethätigt, ſo erfolgt die Entſcheidung
      bezüglich der Annahme doch erſt, nachdem
      ich ihm das fertige Werk\pwindex{deutsche Lyrik des 19. Jahrhunderts. Eine poetische Revue@\emph{Die deutsche Lyrik des 19. Jahrhunderts. Eine poetische Revue}|pwv} vorgelegt habe.
            Ich kann alſo für das thatſächliche Erſcheinen des Werkes\pwindex{deutsche Lyrik des 19. Jahrhunderts. Eine poetische Revue@\emph{Die deutsche Lyrik des 19. Jahrhunderts. Eine poetische Revue}|pwv} keineswegs garantiren.\pend
           
\pstart
           Im Falle Sie auf meinen Vorſchlag eingehen ſollten, bitte ich auch um Angabe
      von Geburtsdatum und Ort ſowie der
      (eventuellen) Titel \uline{lyrischer} Werke aus
      Ihrer Feder.\pend
           
\pstart
           Indem ich hoffe, Ihnen mit meiner Bitte nicht läſtig zu fallen, zeichne ich, Herr Doktor, {\\[\baselineskip]}Hochachtungsvoll\spacefill\mbox{Theodor vSosnosky}\pend
           \leftskip=0em{}
\pstart
           \noindent{}27/12 99.\pend
           
\pstart
           dzt. \textsc{Kremsmünster}, \textsc{Ober-Oesterreich}\oindex{Kremsmünster@\textbf{Kremsmünster}, \emph{Verwaltungsgebiet}|pw}\pend
           \selectlanguage{ngerman}\endnumbering\briefempfaengerindex{Schnitzler, Arthur@\textsc{Schnitzler, Arthur}!zzzSosnosky, Theodor von@\emph{von Theodor von Sosnosky}!1899-12-271@{27. 12. 1899}|)be}\mylabel{L04225h}
\begin{anhang}
\end{anhang}\newcommand{\dateiname}{L04225}\newcommand{\titel}{Theodor von Sosnosky an Arthur Schnitzler, 27. 12. 1899}\newcommand{\editorInnen}{Herausgegeben von Jahnke, SelmaMüller, Martin Anton}\input{../tex-inputs/latex-leseansicht-abspann}
